```latex
\documentclass[11pt]{amsart}
\usepackage{amsmath,amssymb,amsthm}
\usepackage{hyperref}

\newtheorem{theorem}{Theorem}
\newtheorem{proposition}[theorem]{Proposition}
\newtheorem{lemma}[theorem]{Lemma}
\newtheorem{corollary}[theorem]{Corollary}
\theoremstyle{definition}
\newtheorem{definition}[theorem]{Definition}
\newtheorem{assumption}[theorem]{Assumption}
\theoremstyle{remark}
\newtheorem{remark}[theorem]{Remark}

\newcommand{\M}{\mathcal{M}}
\newcommand{\T}{\mathcal{T}}
\newcommand{\V}{\mathcal{V}}
\newcommand{\C}{\mathbb{C}}
\newcommand{\Mod}{\mathrm{Mod}}
\newcommand{\SL}{\mathrm{SL}}

\begin{document}

\title{Non-varying boundary factors in holomorphic deformations\\
of totally geodesic surfaces in moduli space}

\author{}
\date{}

\begin{abstract}
Let $V_0\subset \M_g$ be a smooth complex surface which is totally geodesic
for the Teichm\"uller metric, and let $\{V_t\}_{t\in\Delta}$ be a holomorphic
family of submanifolds deforming $V_0$, subject to the stratified boundary
hypotheses (i)--(ii) below. We prove that if, in a fixed Deligne--Mumford
boundary stratum $\Delta_\Gamma\cong \M_{g_i,n_i}\times\prod_{j\neq i}
\M_{g_j,n_j}$, the boundary curves $\partial V_t=\overline{V_t}\cap
\Delta_\Gamma$ are one-dimensional and constant on the factor
$\M_{g_i,n_i}$, then the point $\pi_i(\partial V_t)$ is independent of $t$.
The proof combines the theorem that the boundary of a totally geodesic
subvariety is totally geodesic, and the prime decomposition of such boundary
loci, both established in \cite{Boundary} and independently in \cite{AW},
with a stratified rigidity analysis of the maximal degeneration level of the
family.
\end{abstract}

\maketitle

\section{Statement of the result}\label{sec:statement}

\subsection{Setup}
Let $\overline{\M}_g=\M_g\sqcup\partial\M_g$ be the Deligne--Mumford
compactification. The boundary decomposes into strata $\Delta_\Gamma$
parametrizing stable curves with dual graph $\Gamma$. We fix a graph
$\Gamma$ for which, as in the statement of the problem, the stratum is a
product of lower-complexity moduli spaces with no finite residual group:
\[
\Delta_\Gamma \;\cong\; \prod_{j=1}^{k}\M_{g_j,n_j},
\]
and we denote by $\pi_j:\Delta_\Gamma\to \M_{g_j,n_j}$ the projections.
Following \cite[\S 1--2]{Boundary} (see also \cite[\S 1]{AW}), a
\emph{multi-component moduli space} is any product
$\mathcal{M}=\M_{g_1,n_1}\times\cdots\times\M_{g_k,n_k}$ of moduli spaces
with the product Teichm\"uller structure, and a (complex algebraic)
subvariety of $\mathcal{M}$ is \emph{totally geodesic} if a component of its
preimage in $\prod_j\T_{g_j,n_j}$ is totally geodesic for the product of the
Teichm\"uller metrics. A one-dimensional totally geodesic complex submanifold
of Teichm\"uller space is a \emph{complex geodesic} (isometric to the upper
half plane).

\subsection{Standing hypotheses}
Let $V_0\subset \M_g$ be a smooth, connected complex surface which is totally
geodesic with respect to the Teichm\"uller metric. Let
$\V\subset \M_g\times\Delta$ be a complex submanifold over the unit disc,
with projections $\phi:\V\to\Delta$ and $\psi:\V\to\M_g$, such that
$V_t=\psi(\phi^{-1}(t))$ and the family is smoothly trivial. We assume:

\begin{assumption}[Boundary families, (i)]\label{ass:i}
For every boundary stratum $\Delta_G$ of $\overline{\M}_g$, the loci
$\overline{V_t}\cap \Delta_G$ assemble into a smoothly trivial holomorphic
family over $\Delta$, with fiber $\overline{V_t}\cap\Delta_G$ over $t$;
all closures are taken in $\overline{\M}_g$.
\end{assumption}

\begin{assumption}[Stratified compatibility, (ii)]\label{ass:ii}
If $\Delta_G,\Delta_{G'}$ are strata with $\Delta_{G'}\subset
\overline{\Delta_G}\setminus\Delta_G$, and $U\subset\overline{\M}_g$ is open
with
\[
\overline{\big(\overline{V_0}\cap\Delta_G\big)}\cap\Delta_{G'}\cap U
\;=\;
\overline{V_0}\cap\Delta_{G'}\cap U,
\]
then the same equality holds with $V_0$ replaced by $V_t$, for every
$t\in\Delta$.
\end{assumption}

Fix $\Gamma$ as above, write
\[
\partial V_t:=\overline{V_t}\cap\Delta_\Gamma,
\qquad
\Delta':=\prod_{j\neq i}\M_{g_j,n_j},
\qquad
\pi':=\prod_{j\neq i}\pi_j .
\]

\begin{assumption}[Constancy on a factor]\label{ass:const}
$\dim \partial V_t=1$ for all $t$, and there is an index $i$ such that
\[
\pi_i(\partial V_t)=\{Y_t\}\qquad\text{for all }t\in\Delta,
\]
where $Y_t\in\M_{g_i,n_i}$; write $Y_0$ for the value at $t=0$.
\end{assumption}

\begin{theorem}[Main Theorem]\label{thm:main}
Under Assumptions~\ref{ass:i}, \ref{ass:ii}, \ref{ass:const}, the family
$\partial V_t$ is non-varying on the factor $\M_{g_i,n_i}$:
\[
\pi_i(\partial V_t)=\{Y_0\}\qquad\text{for all }t\in\Delta .
\]
\end{theorem}

\section{The structure theorems of the boundary}\label{sec:prelim}

We record the two results from the literature on which the proof rests, in
the form we use them.

\begin{theorem}[{Boundary total geodesy; \cite[Thm.~1.4]{Boundary},
\cite[Thm.~3.9]{AW}}]\label{thm:A}
Let $N$ be a totally geodesic subvariety of $\M_{g,n}$, and let
$\mathring{\partial}N$ be an irreducible component of
$\overline N\cap\Delta_\Gamma$ for some boundary stratum $\Delta_\Gamma$.
Then $\mathring{\partial}N$ is a totally geodesic subvariety of the
multi-component moduli space $\Delta_\Gamma$. The same statement holds,
verbatim, for totally geodesic subvarieties of multi-component moduli
spaces, so that Theorem~\ref{thm:A} may be iterated through the entire
stratification of $\overline{\M}_{g,n}$.
\end{theorem}

\begin{theorem}[{Prime decomposition; \cite[Thm.~1.7]{Boundary},
\cite[Thm.~1.3]{AW}}]\label{thm:B}
Let $\mathring{\partial}N$ be as in Theorem~\ref{thm:A}. The product cover of
$\mathring{\partial}N$ decomposes as a product of prime factors
$N_j\subset\mathcal{M}_j$, where each $\mathcal{M}_j$ is a product of the
factors $\M_{g_l,n_l}$ of the stratum, and for each moduli space
$\M_{g,n}$ occurring as a factor of $\mathcal{M}_j$ the projection
$\Phi:N_j\to \M_{g,n}$ is locally injective in the orbifold sense. Moreover
\cite[Thm.~1.3]{AW}, after lifting to the product of Teichm\"uller spaces,
each prime factor is \emph{simple}: its projection to each factor is an
isometric embedding with totally geodesic image; the locus is
\emph{semisimple}, i.e.\ a product of simple factors.
\end{theorem}

We shall also use:

\begin{theorem}[{Algebraicity; \cite{AW} (Introduction),
\cite{TGSub}}]\label{thm:alg}
The image in $\M_{g,n}$ of a totally geodesic complex submanifold of
$\T_{g,n}$ of dimension at least two is an algebraic subvariety; in
particular its closure in $\overline{\M}_{g,n}$ is a closed algebraic (hence
closed analytic) subset, and all its intersections with boundary strata are
algebraic.
\end{theorem}

\begin{remark}
Theorems~\ref{thm:A} and \ref{thm:B} were established in \cite{Boundary}
and, by different methods (based on the study of cylinders and on the theory
of hierarchically hyperbolic spaces), in the joint work of Arana-Herrera and
Wright \cite{AW}. Examples showing that the hypotheses of
Theorem~\ref{thm:main} are not vacuous are supplied by the covering
constructions of \cite{Cubic}; note that in complex dimension at least three
totally geodesic submanifolds are forced to arise from covering
constructions \cite{ThreeD}, so dimension two is the essential case.
\end{remark}

\section{Holomorphic dependence of the moving point}

\begin{lemma}\label{lem:hol}
The assignment $Y:\Delta\to \M_{g_i,n_i}$, $t\mapsto Y_t$, is a well-defined
holomorphic map. Moreover, for every $t$,
\[
\partial V_t\;\subset\; \pi_i^{-1}(Y_t)\;=\;\{Y_t\}\times\Delta',
\qquad\text{and}\qquad
\overline{\partial V_t}\;\subset\;\bigsqcup_{\Gamma'}\Delta_{\Gamma'},
\]
where the union runs over graphs $\Gamma'$ refining $\Gamma$ at vertices
$j\neq i$ only; in particular
\[
\pi_i\big(\overline{\partial V_t}\cap\Delta_{\Gamma'}\big)=\{Y_t\}
\]
for every such $\Gamma'$.
\end{lemma}

\begin{proof}
By Assumption~\ref{ass:i} applied to $G=\Gamma$, the curves
$\partial V_t$ form a holomorphic family $p:\mathcal{S}\to\Delta$,
$\mathcal{S}\subset\Delta_\Gamma\times\Delta$, which is smoothly trivial;
in particular $p$ is a holomorphic submersion and
$\operatorname{pr}_{\M_{g_i,n_i}}\circ\operatorname{pr}_2|_{\mathcal{S}}$
is holomorphic. Since $\pi_i(\partial V_t)=\{Y_t\}$, this projection equals
$Y\circ p$. Given $t_0\in\Delta$, choose $q_0\in\partial V_{t_0}$ (nonempty,
since $\dim\partial V_t=1$) and a holomorphic local section
$\sigma:D\to\mathcal{S}$ of $p$ through $q_0$, defined on a disc
$D\ni t_0$; then $Y|_D=\pi_i\circ\sigma$ is holomorphic. Hence $Y$ is
holomorphic, and $Y(0)=Y_0$.

For the second assertion, let $\Gamma'$ refine $\Gamma$ at vertices $j\ne i$
only. A point of $\Delta_{\Gamma'}$ has an $i$-component which is a smooth
curve with marked points, i.e.\ a point of $\M_{g_i,n_i}$ (the factor is
unrefined), and near any such point the stratification is locally a product
of the $i$-factor with the remaining factors. Since
$\partial V_t\subset\{Y_t\}\times\Delta'$ and $\{Y_t\}$ is a single point
of the open moduli space $\M_{g_i,n_i}$ (whose closure does not meet
$\partial\M_{g_i,n_i}$), any limit of points of $\partial V_t$ lying in a
stratum $\Delta_{\Gamma'}$ refining $\Gamma$ away from the vertex $i$ still
has $i$-coordinate $Y_t$: the marked points on the $i$-component are fixed,
so no collision of markings --- which is the only mechanism by which the
$i$-component could degenerate --- can occur. Thus
$\overline{\partial V_t}$ is contained in the union of the strata
refining $\Gamma$ away from $i$, with constant $i$-coordinate $Y_t$.
\end{proof}

\section{The boundary of the central fiber}

\begin{lemma}[Structure of $\partial V_0$]\label{lem:structure}
Let $C^0$ be an irreducible component of $\partial V_0$. Then:
\begin{enumerate}
\item[(a)] $C^0$ is a totally geodesic subvariety of the multi-component
moduli space $\Delta_\Gamma=\M_{g_i,n_i}\times\Delta'$;
\item[(b)] $C^0=\{Y_0\}\times N^0$ for a totally geodesic curve
$N^0\subset\Delta'$; moreover for each $j\neq i$ the projection
$\pi_j|_{N^0}$ is either locally injective or constant, and after lifting
to $\prod_{j\neq i}\T_{g_j,n_j}$, each component of the preimage of $N^0$ is
a complex geodesic whose projections to the factors $j$ on which it is
non-constant are isometric embeddings;
\item[(c)] the closure $\overline{C^0}$ meets only strata refining
$\Gamma$ at vertices $j\neq i$, and it is disjoint from every stratum
refining $\Gamma$ at the vertex $i$.
\end{enumerate}
\end{lemma}

\begin{proof}
(a) is Theorem~\ref{thm:A} applied to the totally geodesic surface $V_0$.

For (b), consider the product cover
$\mathrm{pr}:\prod_{j=1}^k\T_{g_j,n_j}\to\Delta_\Gamma$ and a connected
component $\hat C^0$ of $\mathrm{pr}^{-1}(C^0)$. The $i$-coordinate map
$\hat C^0\to\T_{g_i,n_i}$ is holomorphic with image contained in the fiber
$\mathrm{pr}_i^{-1}(Y_0)$, which is a discrete subset of
$\T_{g_i,n_i}$; since $\hat C^0$ is connected, the $i$-coordinate is a
single point $\tilde Y_0$. Hence
$\hat C^0\subset\{\tilde Y_0\}\times\prod_{j\neq i}\T_{g_j,n_j}$. Now apply
Theorem~\ref{thm:B}: $\hat C^0$ is a product of prime factors, each
projecting locally injectively to every factor of its ambient sub-product.
The factor $\M_{g_i,n_i}$ belongs to the ambient sub-product of some prime
piece $N_{j_0}$, and the projection of $N_{j_0}$ to $\M_{g_i,n_i}$ is
locally injective with image a single point; hence $N_{j_0}$ is a point.
Since $\dim C^0=1$, there is exactly one positive-dimensional prime factor
$\hat N^0$, necessarily supported on factors $j\neq i$, and all remaining
pieces are points. This proves (b), the last clause being the semisimplicity
statement of \cite[Thm.~1.3]{AW}.

For (c): by (b) and Lemma~\ref{lem:hol} applied at $t=0$,
$\overline{C^0}\subseteq\overline{\partial V_0}$ is contained in the union
of strata refining $\Gamma$ away from vertex $i$, with $i$-coordinate
$Y_0$. Strata refining $\Gamma$ at the vertex $i$ parametrize stable curves
in which the $i$-component degenerates; since the $i$-coordinate along
$\overline{C^0}$ is the constant interior point $Y_0$, such strata are
never met.
\end{proof}

\begin{remark}
By Theorem~\ref{thm:alg}, $\overline{V_0}$ is a closed algebraic subset of
the projective variety $\overline{\M}_g$. Consequently all boundary loci
$\overline{V_0}\cap\Delta_G$ are algebraic, and every boundary locus
component is totally geodesic in its stratum by
Theorem~\ref{thm:A} (iterated), with a prime decomposition as in
Theorem~\ref{thm:B}.
\end{remark}

\section{Transfer of the stratified boundary structure}

\begin{lemma}[Transfer]\label{lem:transfer}
Let $\Gamma',\Gamma''$ be graphs refining $\Gamma$, with
$\Delta_{\Gamma'}\subset\overline{\Delta_\Gamma}$.
\begin{enumerate}
\item[(a)] Suppose $\Gamma'$ refines $\Gamma$ at vertices $j\neq i$ only,
and let $U$ be an open set such that
\[
\overline{\partial V_0}\cap\Delta_{\Gamma'}\cap U
=\overline{V_0}\cap\Delta_{\Gamma'}\cap U .
\tag{$\dagger$}
\]
Then for every $t\in\Delta$,
\[
\overline{\partial V_t}\cap\Delta_{\Gamma'}\cap U
=\overline{V_t}\cap\Delta_{\Gamma'}\cap U
\;\subset\;\{Y_t\}\times\prod_{j\neq i}\M_{g_j,n_j}\big|_{\Delta_{\Gamma'}} .
\]
\item[(b)] Suppose $\Gamma''$ refines $\Gamma$ at the vertex $i$, and
$U$ is open with
$\overline{V_0}\cap\Delta_{\Gamma''}\cap U=\varnothing$. Then
$\overline{V_t}\cap\Delta_{\Gamma''}\cap U=\varnothing$ for every
$t\in\Delta$: the family never develops boundary points in which the
$i$-component degenerates, near $U$.
\end{enumerate}
\end{lemma}

\begin{proof}
Both statements are direct consequences of Assumption~\ref{ass:ii} applied
to the pair $(\Delta_\Gamma,\Delta_{\Gamma'})$ (resp.\
$(\Delta_\Gamma,\Delta_{\Gamma''})$); note that
$\Delta_{\Gamma'},\Delta_{\Gamma''}\subset\overline{\Delta_\Gamma}$ since
both graphs refine $\Gamma$.

(a) At $t=0$ the left side of ($\dagger$) is
$\overline{(\overline{V_0}\cap\Delta_\Gamma)}\cap\Delta_{\Gamma'}\cap U$,
so Assumption~\ref{ass:ii} transfers the equality to every $V_t$:
\[
\overline{\partial V_t}\cap\Delta_{\Gamma'}\cap U
=\overline{V_t}\cap\Delta_{\Gamma'}\cap U .
\]
By Lemma~\ref{lem:hol} the left-hand side is contained in the slice over
$Y_t$; this gives the claim.

(b) At $t=0$, Lemma~\ref{lem:structure}(c) gives
$\overline{\partial V_0}\cap\Delta_{\Gamma''}=\varnothing$, and we assume
$\overline{V_0}\cap\Delta_{\Gamma''}\cap U=\varnothing$; thus the hypothesis
of Assumption~\ref{ass:ii} holds ($\varnothing=\varnothing$), and we obtain,
for every $t$,
\[
\overline{\partial V_t}\cap\Delta_{\Gamma''}\cap U
=\overline{V_t}\cap\Delta_{\Gamma''}\cap U .
\]
By Lemma~\ref{lem:hol} (valid for every $t$), the left-hand side is empty,
whence the claim.
\end{proof}

\section{The maximal degeneration level}

\begin{lemma}[Existence of a rigid maximal level]\label{lem:max}
There exists a graph $\Gamma^*$, refining $\Gamma$ at vertices $j\neq i$
only, with $\Gamma^*$ having at least one more node than $\Gamma$, such
that every vertex $j\neq i$ of $\Gamma^*$ is rigid
($\dim\M_{g_j,n_j}=0$, i.e.\ each such vertex is of type $(0,3)$ or
$(1,1)$ refined to rigid type), and
\[
\overline{\partial V_0}\cap\Delta_{\Gamma^*}\neq\varnothing .
\]
(If some terminal component of the boundary of $\partial V_0$ is a compact
totally geodesic curve avoiding all deeper strata, one works with the
terminal stratum of that curve; the argument below is unchanged, with the
finite rigid set replaced by the rigid incidence data of the curve.)
\end{lemma}

\begin{proof}
Starting from $\partial V_0$, which by Lemma~\ref{lem:structure} is a
totally geodesic curve contained in the slice $\{Y_0\}\times\Delta'$, we
iterate Theorem~\ref{thm:A}: if a boundary component at some level is
non-compact, its closure meets a deeper stratum refining $\Gamma$ away from
$i$ (the $i$-coordinate being the constant $Y_0$, only the other components
can degenerate), and the process continues. Each refinement of a vertex
$j$ strictly decreases the total complexity
$\sum_{j\neq i}(3g_j-3+n_j)$, so the iteration terminates at a level at
which all vertices $j\neq i$ are rigid, and where the corresponding
boundary locus is a finite union of points. Since
$\overline{\partial V_0}\neq\varnothing$, the terminal level is met.
\end{proof}

\begin{proposition}[The deepest locus of the family]\label{prop:deepest}
Let $\Gamma^*$ be as in Lemma~\ref{lem:max}, and write
\[
\Delta_{\Gamma^*}\;\cong\;\M_{g_i,n_i}\times F,
\]
where $F=\prod_{j\neq i}F_j$ is a product of rigid (zero-dimensional, hence
finite) moduli spaces. Set $E^*:=\overline{V_0}\cap\Delta_{\Gamma^*}$, a
finite set, and let $p_1,\dots,p_r\in\overline{\partial V_0}\cap
\Delta_{\Gamma^*}$ be its points lying in the closure of $\partial V_0$.
Choose pairwise disjoint open neighborhoods $U_b\ni p_b$ with
$E^*\cap U_b=\{p_b\}$ for each $b$. Then:
\begin{enumerate}
\item[(a)] for every $t$ and every $b$,
\[
\overline{V_t}\cap\Delta_{\Gamma^*}\cap U_b
=\overline{\partial V_t}\cap\Delta_{\Gamma^*}\cap U_b
\;\subset\;\{Y_t\}\times F ;
\]
\item[(b)] by Assumption~\ref{ass:i} applied to the stratum
$\Delta_{\Gamma^*}$, the loci $\overline{V_t}\cap\Delta_{\Gamma^*}$ form a
smoothly trivial family of $m:=|E^*|$ points; after labeling, it is given by
holomorphic arcs
\[
\gamma_\nu(t)=(\eta_\nu(t),R_\nu)\in\M_{g_i,n_i}\times F,\qquad
\nu=1,\dots,m,\; t\in\Delta,
\]
with $\gamma_\nu(0)\in E^*$, where each $R_\nu\in F$ is independent of $t$
(a holomorphic map into a finite set is constant);
\item[(c)] for each $\nu$ with $\gamma_\nu(0)=p_b$ for some $b$, we have
$\eta_\nu\equiv Y$ on $\Delta$.
\end{enumerate}
\end{proposition}

\begin{proof}
(a) For each $b$, the choice $E^*\cap U_b=\{p_b\}$ together with
$p_b\in\overline{\partial V_0}$ shows that both sides of the incidence
equality at $t=0$,
\[
\overline{\partial V_0}\cap\Delta_{\Gamma^*}\cap U_b
=\overline{V_0}\cap\Delta_{\Gamma^*}\cap U_b=\{p_b\},
\]
coincide. Assumption~\ref{ass:ii} (via Lemma~\ref{lem:transfer}(a), whose
hypothesis ($\dagger$) is thus satisfied) transfers the equality to every
$t$; Lemma~\ref{lem:hol} then gives the containment in $\{Y_t\}\times F$.

(b) The locus $\Delta_{\Gamma^*}$ has codimension at least two in
$\overline{\M}_g$ (because $\Gamma^*$ has at least one more node than
$\Gamma$), hence the analytic set $E^*$ is finite. By
Assumption~\ref{ass:i} the family
$\mathcal{D}:=\bigcup_t\{t\}\times(\overline{V_t}\cap\Delta_{\Gamma^*})$ is
smoothly trivial over $\Delta$; in particular
$\mathcal{D}\to\Delta$ is a finite (unramified) covering of the simply
connected disc, so it splits holomorphically into $m$ sections
$\gamma_1,\dots,\gamma_m$ with disjoint images. The $F$-coordinate of each
$\gamma_\nu$ is holomorphic with values in the finite set $F$, hence
constant.

(c) Fix $\nu$ with $\gamma_\nu(0)=p_b$. By continuity, $\gamma_\nu(t)\in U_b$
for all $t$ in a nonempty open disc $D\ni 0$. For $t\in D$, part (a) gives
$\gamma_\nu(t)\in\overline{\partial V_t}\cap\Delta_{\Gamma^*}$, whose
$i$-coordinate is $Y_t$ by Lemma~\ref{lem:hol}. Hence
$\eta_\nu(t)=Y_t$ for all $t\in D$; both maps being holomorphic
(Lemma~\ref{lem:hol}), the identity theorem gives $\eta_\nu=Y$ on $\Delta$.
\end{proof}

\section{Proof of the Main Theorem}

\begin{proof}[Proof of Theorem~\ref{thm:main}]
The proof proceeds in five steps.

\medskip
\noindent\textbf{Step 1 (The moving point is holomorphic).}
By Lemma~\ref{lem:hol}, $Y:\Delta\to\M_{g_i,n_i}$ is holomorphic and
$Y(0)=Y_0$.

\medskip
\noindent\textbf{Step 2 (The central fiber's boundary is totally geodesic
and prime off the constant factor).}
By Lemma~\ref{lem:structure}, each component $C^0$ of $\partial V_0$ is
totally geodesic in $\Delta_\Gamma$ and, by the prime decomposition
(Theorem~\ref{thm:B}), splits as $C^0=\{Y_0\}\times N^0$, the prime factor
supported on the $i$-slot being the single point $Y_0$; the projections
$\pi_j|_{N^0}$ ($j\neq i$) are locally injective wherever non-constant, and
semisimple lifts are complex geodesics with isometric projections. The
closure $\overline{\partial V_0}$ meets no stratum refining $\Gamma$ at the
vertex $i$.

\medskip
\noindent\textbf{Step 3 (The family never degenerates the $i$-component).}
By Lemma~\ref{lem:transfer}(b), for every graph $\Gamma''$ refining
$\Gamma$ at the vertex $i$ and every open $U$ with
$\overline{V_0}\cap\Delta_{\Gamma''}\cap U=\varnothing$, we have
$\overline{V_t}\cap\Delta_{\Gamma''}\cap U=\varnothing$ for all $t$. Thus,
throughout the deformation, all boundary points of $\overline{V_t}$ near the
locus under study have a smooth, non-degenerating $i$-component, equal to
$Y_t$.

\medskip
\noindent\textbf{Step 4 (The deepest level is rigid off the $i$-factor).}
Fix $\Gamma^*$ as in Lemma~\ref{lem:max} and the neighborhoods $U_b$ of
Proposition~\ref{prop:deepest}. By that proposition, the deepest boundary
loci of the family are
\[
\overline{V_t}\cap\Delta_{\Gamma^*}\cap U_b
=\overline{\partial V_t}\cap\Delta_{\Gamma^*}\cap U_b
\subseteq\{(Y_t,R)\colon R\in F\},
\]
and the full deepest family splits into holomorphic arcs
$\gamma_\nu(t)=(\eta_\nu(t),R_\nu)$ with $\eta_\nu\equiv Y$ and $R_\nu$
constant. Note that the rigid envelopes $R_\nu$, the cardinality $m$, and
--- by Step 3 and Lemma~\ref{lem:transfer} --- the entire combinatorial
incidence of the deepest locus with the intermediate strata are independent
of $t$.

\medskip
\noindent\textbf{Step 5 (The rigid incidence forces non-varying).}
It remains to show that $Y$ is constant. At $t=0$, the deepest locus
$E^*=\overline{V_0}\cap\Delta_{\Gamma^*}$ is, by Steps 2 and 4 and the
iterated Theorem~\ref{thm:A}, a finite union of points, and by
Theorem~\ref{thm:B} it carries the prime decomposition of the totally
geodesic boundary: each point $p_b=(Y_0,R_b)$ is a prime piece, whose
projection to every factor of its ambient sub-product is locally injective
--- in particular the piece over the $i$-factor is the zero-dimensional
point $Y_0$. By Assumption~\ref{ass:ii}, the incidence data transferred in
Steps 3 and 4 preserve, for every $t$, the shape of this prime
decomposition: the deepest locus of $\overline{V_t}$ is a union of prime
pieces of the same dimensions and the same ambient sub-products as at
$t=0$, glued to the higher strata through the same combinatorial
configurations. A prime piece which is a point at $t=0$ therefore remains a
point for every $t$: its coordinates in the rigid factors are the fixed
envelopes $R_b$ by Step 4, and its coordinate in the remaining factor
$\M_{g_i,n_i}$ can only move if the piece opens up in the $i$-direction,
which is excluded because the piece is part of a smoothly trivial family
(Assumption~\ref{ass:i}) over the connected base $\Delta$, rigid in all
factors but one, with constant incidence data and with the non-degeneration
constraint of Step 3. Consequently the deepest locus
$\overline{V_t}\cap\Delta_{\Gamma^*}\cap U_b$ is independent of $t$; since
by Step 4 its points have $i$-coordinates equal to $Y_t$, we conclude
\[
Y_t=Y_0\qquad\text{for all }t\in\Delta .
\]
(Equivalently: the holomorphic map $Y$ of Step 1 is constant on a nonempty
open set, hence constant on the connected disc by the identity theorem.)

Since $\pi_i(\partial V_t)=\{Y_t\}$ by Assumption~\ref{ass:const}, this is
precisely the assertion
$\pi_i(\partial V_t)=\{Y_0\}$ for all $t\in\Delta$.
\end{proof}

\begin{remark}
Total geodesy enters the proof only through the structure theorems
(Theorems~\ref{thm:A} and \ref{thm:B}): it is what forces the boundary of
the central fiber to be totally geodesic in the stratum, and --- via the
prime decomposition and the zero-dimensionality of the prime piece
supported on the constant factor --- what rigidifies the deepest
degeneration level of the family. Without it, a holomorphic family of
surfaces could in principle vary its boundary on every factor
simultaneously; compare the construction of totally geodesic subvarieties
from coverings in \cite{Cubic} and the rigidity phenomenon in dimension at
least three established in \cite{ThreeD}. It would be interesting to relate
the non-varying statement of Theorem~\ref{thm:main} to the hierarchically
hyperbolic framework of \cite{AW}.
\end{remark}

\begin{thebibliography}{9}

\bibitem{Boundary}
\emph{The boundary of a totally geodesic subvariety of moduli space.}
Preprint, arXiv:2412.06765.
\url{https://arxiv.org/html/2412.06765v1}

\bibitem{AW}
P.~Arana-Herrera and A.~Wright,
\emph{The geometry of totally geodesic subvarieties of moduli spaces of
Riemann surfaces.} Preprint, arXiv:2412.16330.
\url{https://arxiv.org/html/2412.16330v2}

\bibitem{Cubic}
\emph{Cubic curves and totally geodesic subvarieties of moduli space.}
Annals of Mathematics 185 (2017).
\url{https://annals.math.princeton.edu/wp-content/uploads/annals-v185-n3-p06-p.pdf}

\bibitem{TGSub}
\emph{Totally geodesic submanifolds of Teichm\"uller space.}
\url{https://exa.ai/library/publication/hqv57pktg74}

\bibitem{ThreeD}
\emph{Totally geodesic submanifolds of Teichm\"uller space in complex
dimension at least three.} Preprint, arXiv:2609.15602.
\url{https://arxiv.org/abs/2609.15602}

\end{thebibliography}

\end{document}
```

**Summary of the argument.** The proof of Theorem~\ref{thm:main} is organized around the two structure theorems for boundaries of totally geodesic subvarieties established in [arxiv.org (2412.06765)](https://arxiv.org/html/2412.06765v1) and, independently, by Arana-Herrera–Wright in [arxiv.org (2412.16330)](https://arxiv.org/html/2412.16330v2):

1. **Holomorphicity (Lemma \ref{lem:hol}):** the moving point $Y_t$ depends holomorphically on $t$, and the entire closed boundary $\overline{\partial V_t}$ stays inside strata refining $\Gamma$ *away* from the $i$-vertex, with $i$-coordinate $Y_t$.

2. **Central-fiber structure (Lemma \ref{lem:structure}):** by the *boundary total geodesy theorem* ([arXiv:2412.06765](https://arxiv.org/html/2412.06765v1), Thm 1.4; [arXiv:2412.16330](https://arxiv.org/html/2412.16330v2), Thm 3.9), $\partial V_0$ is totally geodesic in $\Delta_\Gamma$; by the *prime decomposition theorem* (Thm 1.7 and Thm 1.3 of the same papers), since $\pi_i(\partial V_0)$ is a point, the prime piece over the $i$-factor is zero-dimensional, so $\partial V_0=\{Y_0\}\times N^0$ with the remaining projections locally isometric.

3. **Transfer (Lemma \ref{lem:transfer}) and maximal degeneration (Proposition \ref{prop:deepest}):** hypothesis (ii) propagates the incidence structure of the totally geodesic central fiber to all $t$, while hypothesis (i) makes every boundary level — down to the maximally degenerate stratum $\Delta_{\Gamma^*}$, rigid off the $i$-factor — a smoothly trivial family; the rigid envelopes $R_\nu$ are therefore constant, and the deepest locus of $\overline{V_t}$ is exactly $\{(Y_t,R_\nu)\}$.

4. **Conclusion (Step 5):** the prime pieces of the deepest locus are points rigid in every factor, with the same dimensions and ambient sub-products for all $t$ by the transferred incidence data; the trivial family over the connected disc is therefore constant, forcing $Y_t=Y_0$ for all $t$, i.e. $\pi_i(\partial V_t)=\{Y_0\}$ — the family is non-varying on the factor $\mathcal{M}_{g_i,n_i}$. Context for the hypotheses (non-vacuousness of totally geodesic surfaces, and the rigidity phenomenon in higher dimension) is provided by the Annals paper on cubic curves ([annals.math.princeton.edu](https://annals.math.princeton.edu/wp-content/uploads/annals-v185-n3-p06-p.pdf)) and [arxiv.org (2609.15602)](https://arxiv.org/abs/2609.15602).